% Источник решений: «4 Семинар Ивана.pdf», страницы 2--6.
% Условия: conditions/sem_04.pdf.
\section{Независимость и дискретные распределения}

\noindent Условия:
\href{run:conditions/sem_04.pdf}{семинар 4 (PDF)}.

\subsection{Задача 1}

\begin{solution}
    \textbf{Независимость событий.}

    \textbf{а)} Дано: $\mathcal{A}\perp\!\!\!\perp\mathcal{B}$.
    % В рукописи справа пропущено вычитание A.
    \[
        \overline{\mathcal{A}}\cap\mathcal{B}
        =(\Omega\setminus\mathcal{A})\cap\mathcal{B}
        =\mathcal{B}\setminus\mathcal{A}.
    \]
    \begin{center}
        \begin{tikzpicture}[scale=0.8]
            \draw (0,0) ellipse (2.5 and 1.3);
            \node at (-1.9,0.45) {$\Omega$};
            \begin{scope}
                \clip (0.5,0) circle (0.8);
                \fill[blue!20] (-2,-1) rectangle (2,1);
                \fill[solution_back] (-0.5,0) circle (0.8);
            \end{scope}
            \draw (-0.5,0) circle (0.8);
            \draw (0.5,0) circle (0.8);
            \node at (-0.9,0.4) {$\mathcal{A}$};
            \node at (0.9,0.4) {$\mathcal{B}$};
        \end{tikzpicture}
    \end{center}
    \[
        \begin{aligned}
            \Pr(\overline{\mathcal{A}}\cap\mathcal{B})
            &=\Pr(\mathcal{B})-\Pr(\mathcal{A}\cap\mathcal{B})\\
            &=\Pr(\mathcal{B})-\Pr(\mathcal{A})\Pr(\mathcal{B})\\
            &=\Pr(\mathcal{B})(1-\Pr(\mathcal{A}))
             =\Pr(\mathcal{B})\Pr(\overline{\mathcal{A}}).
        \end{aligned}
    \]
    Следовательно, $\overline{\mathcal{A}}\perp\!\!\!\perp\mathcal{B}$.

    \medskip
    \textbf{Модель: неклассическая дискретная. Схема Бернулли.}
    \[
        \Omega=\{0,1\}^n,\qquad
        \omega=(i_1,\ldots,i_n),\qquad |\Omega|=2^n,\qquad n\geq2.
    \]
    \[
        \Pr(\{\omega\})
        =p^{\sum_{j=1}^n i_j}(1-p)^{n-\sum_{j=1}^n i_j}.
    \]
    \[
        \mathcal{A}=\{\omega:i_1=1\},\qquad
        \mathcal{B}=\{\omega:i_n=1\}.
    \]
    % В рукописи не дописана степень при (1-p).
    \[
        \begin{aligned}
            \Pr(\mathcal{A})
            &=\sum_{\omega:\,i_1=1}\Pr(\{\omega\})\\
            &=p\sum_{i_2,\ldots,i_n\in\{0,1\}}
                p^{\sum_{j=2}^n i_j}
                (1-p)^{n-1-\sum_{j=2}^n i_j}
             =p\cdot1=p.
        \end{aligned}
    \]
    Аналогично $\Pr(\mathcal{B})=p$.
    \[
        \begin{aligned}
            \Pr(\mathcal{A}\cap\mathcal{B})
            &=\sum_{\omega:\,i_1=i_n=1}\Pr(\{\omega\})\\
            &=p^2\sum_{i_2,\ldots,i_{n-1}\in\{0,1\}}
                p^{\sum_{j=2}^{n-1}i_j}
                (1-p)^{n-2-\sum_{j=2}^{n-1}i_j}\\
            &=p^2=\Pr(\mathcal{A})\Pr(\mathcal{B}).
        \end{aligned}
    \]

    \textbf{б)}
    \[
        \begin{aligned}
            \mathcal{A}\perp\!\!\!\perp\mathcal{A}
            &\iff\Pr(\mathcal{A}\cap\mathcal{A})=\Pr(\mathcal{A})^2\\
            &\iff\Pr(\mathcal{A})=\Pr(\mathcal{A})^2\\
            &\iff\Pr(\mathcal{A})=0\ \text{или}\ \Pr(\mathcal{A})=1.
        \end{aligned}
    \]
\end{solution}

\subsection{Задача 2}

\begin{solution}
    \textbf{Модель: дискретная.}

    Дискретное распределение: $\{(x_i,p_i)\}$, где
    \[
        p_i=\Pr\{\xi=x_i\},\qquad \sum_i p_i=1.
    \]
    $\{\xi=x_i\}$ --- событие, элемент $\mathcal{F}$.
    % В рукописи у суммы не указано ограничение на индекс.
    \[
        \Pr_\xi(B)=\sum_{i:\,x_i\in B}p_i.
    \]
    В задаче $x_i\in\{1,2,3\}$:
    \[
        \begin{aligned}
            p_1&=\Pr\{\xi=1\}=\frac13,\\
            p_2&=\Pr\{\xi=2\}=\frac23\cdot\frac12=\frac13,\\
            p_3&=\Pr\{\xi=3\}=\frac23\cdot\frac12\cdot1=\frac13.
        \end{aligned}
    \]
    \[
        \boxed{\left\{\left(1,\frac13\right),
            \left(2,\frac13\right),\left(3,\frac13\right)\right\}}.
    \]
\end{solution}

\subsection{Задача 3}

\begin{solution}
    \textbf{Модель: дискретная. Геометрическое распределение.}
    \[
        \Pr\{\xi=k\}=(1-p)^{k-1}p,\qquad k\in\N,\qquad 0<p\leq1.
    \]
    Здесь $k-1$ неуспехов и затем успех.
    \[
        \eta=\frac{\xi}{2}\bigl(1-(-1)^\xi\bigr),\qquad
        y_i\in\{0,1,3,5,\ldots\}.
    \]
    % Исправлено промежуточное суммирование геометрической прогрессии;
    % итоговая вероятность в рукописи верна.
    \[
        \begin{aligned}
            \Pr\{\eta=0\}
            &=\sum_{n=1}^{\infty}\Pr\{\xi=2n\}
             =p\sum_{n=1}^{\infty}(1-p)^{2n-1}\\
            &=\frac{p(1-p)}{1-(1-p)^2}
             =\frac{p(1-p)}{p(2-p)}
             =\boxed{\frac{1-p}{2-p}}.
        \end{aligned}
    \]
    \[
        \Pr\{\eta=2k+1\}=\Pr\{\xi=2k+1\}
        =\boxed{(1-p)^{2k}p},\qquad k=0,1,2,\ldots
    \]
\end{solution}

\clearpage
\subsection{Задача 4}

\begin{solution}
    \textbf{Модель: дискретная. Распределение Пуассона.}
    \[
        p_k=\Pr\{\xi=k\}=\frac{\lambda^k}{k!}e^{-\lambda},
        \qquad k=0,1,2,\ldots
    \]
    Для дискретных случайных величин
    \[
        \xi\perp\!\!\!\perp\eta
        \iff\forall\,i,j\quad
        \{\omega:\xi(\omega)=x_i\}\perp\!\!\!\perp
        \{\omega:\eta(\omega)=y_j\}.
    \]
    Пусть $\xi\sim\operatorname{Poiss}(\lambda_1)$,
    $\eta\sim\operatorname{Poiss}(\lambda_2)$,
    $\lambda_1,\lambda_2>0$, $\xi\perp\!\!\!\perp\eta$.
    \[
        \zeta=\xi+\eta\in\{0,1,2,\ldots\}.
    \]
    По формуле полной вероятности:
    % В рукописи во второй сумме один раз ошибочно написано i=1.
    \[
        \begin{aligned}
            \Pr\{\zeta=k\}
            &=\sum_{i=0}^{\infty}\Pr\{\xi=i\}
                \Pr\{\xi+\eta=k\mid\xi=i\}\\
            &=\sum_{i=0}^{k}\Pr\{\xi=i\}
                \Pr\{\eta=k-i\mid\xi=i\}\\
            &=\sum_{i=0}^{k}\Pr\{\xi=i\}\Pr\{\eta=k-i\}.
        \end{aligned}
    \]
    Последнее равенство следует из независимости. Тогда
    \[
        \begin{aligned}
            \Pr\{\zeta=k\}
            &=\sum_{i=0}^{k}\frac{\lambda_1^i}{i!}e^{-\lambda_1}
                \frac{\lambda_2^{k-i}}{(k-i)!}e^{-\lambda_2}\\
            &=\frac{e^{-(\lambda_1+\lambda_2)}}{k!}
                \sum_{i=0}^{k}C_k^i\lambda_1^i\lambda_2^{k-i}\\
            &=\frac{e^{-(\lambda_1+\lambda_2)}}{k!}
                (\lambda_1+\lambda_2)^k.
        \end{aligned}
    \]
    \[
        \boxed{\xi+\eta\sim\operatorname{Poiss}(\lambda_1+\lambda_2)}.
    \]
\end{solution}
